% ============================================================
% Subgoal Scale & Reality-Verified Speculative Consumption
% target venue: ICLR 2027 (or NeurIPS/ICML equivalent).
%
% ICLR templates are \usepackage, NOT \documentclass -- the .sty
% is not a .cls. Keep \documentclass{article} + \usepackage{iclr...}.
% Swap iclr2026_conference -> iclr2027_conference the moment that
% folder lands in github.com/ICLR/Master-Template.
%
% DOUBLE-BLIND: keep \iclrfinalcopy commented out and the author
% block anonymized for the entire review period.
% ============================================================
\documentclass{article}
\usepackage{iclr2026_conference,times}
\usepackage{graphicx}
\usepackage{amsmath,amssymb}
\usepackage{booktabs}
\usepackage{multirow}
\usepackage{tabularx}
\usepackage{enumitem}
\usepackage{subcaption}
\usepackage{xcolor}
\usepackage{hyperref}
\usepackage{cleveref}

% \iclrfinalcopy   % <-- leave commented out until accepted

% ---- house macros ----
\newcommand{\ours}[1]{\textbf{#1}}          % highlight our winning cells
\newcommand{\degc}[1]{$#1^\circ$}
\newcommand{\gdm}{\text{GDM}}
% Future-experiment placeholder: renders a red TBD marker.
\newcommand{\TBD}{\textcolor{red}{\textbf{TBD}}}
% A framed red placeholder block standing in for a not-yet-run table.
\newcommand{\tbdblock}[1]{%
  \begin{center}\fcolorbox{red}{red!5}{%
    \parbox{0.9\linewidth}{\centering\color{red}\textbf{[TBD --- #1]}}}%
  \end{center}}

\title{Subgoal Scale and Reality-Verified Speculative Consumption\\
in Hierarchical Latent World-Model Planning}

% ANONYMIZED for double-blind review. Restore the real author block
% only after \iclrfinalcopy is uncommented post-acceptance.
\author{Anonymous authors \\ Paper under double-blind review}

\begin{document}
\maketitle

\begin{center}\fcolorbox{red}{red!10}{\parbox{0.92\linewidth}{\centering
\color{red}\textbf{ARCHIVED DRAFT (v1) --- SUPERSEDED BY \texttt{main\_v2.tex}
(2026-07-13).}\\ \color{red}Kept only as source material for the v2 appendices.
Numbers here predate the derived-config re-stamp (spec cells ran $k{=}8$, not the
derived $k{=}3$; rand-init spec cell was unrun); the ``verifier self-test'' claim in
\S\ref{sec:spec} was RETRACTED 2026-07-13 (it is false rejection --- see v2 \S7.1);
the framing predates the calibration-by-measurement section. Do not cite, compile,
or copy numbers from this file except when porting appendix tables.}}
\end{center}

% ============================================================
\begin{abstract}
We study the FF-JEPA hierarchical latent planner --- a frozen LeWM JEPA encoder, a
diffusion subgoal drafter (``GDM''), and a CEM controller --- on PushT, and identify
\emph{subgoal scale}, the spacing $S$ between drafted subgoals with the control loop
matched to it, as the dominant unexplored lever. Success rate over stride is
non-monotone with a clear interior optimum at $S{=}10$: \ours{$97.2\%/96.4\%$}
(\degc{20}/\degc{5} criteria) versus $91.1\%/86.3\%$ for the paper's native $S{=}25$
under an identical protocol. On the paper's own published protocols the winning
configuration \ours{exceeds the published numbers} at both horizons ($96.1{\to}98.4$
at $t{=}25$, $91.8{\to}97.9$ at $t{=}75$, \degc{20}, $n{=}256$); on random
initialization (its hardest setting, published $82.4\%$) it records \ours{zero
failures in 1024 episodes across four seeds}; and on a fixed-population horizon sweep
the baseline decays monotonically while $S{=}10$ stays flat --- \ours{the gap widens
from $+3.3$ to $+11.9$pp}. A cadence-isolation control shows the mechanism is
\emph{scale matching}, not re-planning frequency. Separately, the one component of
DSpark-style speculative decoding that survives measurement --- a draft--verify--accept
skeleton with the learned verifier replaced by the \emph{achieved environment state}
--- cuts diffusion compute ${\sim}11\times$ versus every-step drafting while remaining
success-rate-neutral at the official criterion. The cost claim decomposes into two factors, stated separately throughout: cheap
drafts ($k{=}50\to4$ at equal-or-better SR --- the inherited 50-step sampler was
convention, not necessity) and verified consumption (${\sim}1.8\times$ fewer calls at
matched $k$, SR-neutral). On a second environment (Reacher, exploratory non-expert
data), the mechanism transfers with zero re-tuning --- SR-neutral at every horizon ---
while drafting itself proves regime-dependent: against the strongest flat baseline at
every replanning rate, drafted subgoals win $+11$ to $+19$pp exactly where the goal
outruns the planning window, and are a measurable detour tax where it does not. Refinement heads, learned confidence
scheduling, blind multi-step commitment, and a deterministic regressor that dominates
the drafter \emph{offline} each measurably \emph{hurt} in closed loop; we report them
as negative results that causally implicate drafter stochasticity.
\end{abstract}

% ============================================================
\section{Introduction}
\label{sec:intro}

Hierarchical latent world models promise long-horizon planning by decomposing a goal
into a chain of reachable subgoals. FF-JEPA \citep{ffjepa} instantiates this on top of
the LeWM JEPA world model \citep{lewm}: a diffusion model drafts future subgoal latents
and a flat CEM controller \citep{cem} steers toward each in turn, re-planning on a fixed
stride. It reports that hierarchy beats flat CEM on long horizons --- but under an
underspecified pipeline (no released code; unstated diffusion configuration; unstated
evaluation population), and with a subgoal stride fixed at a single value.

We take FF-JEPA as a substrate and ask what actually governs its closed-loop success.
Our contributions:

\begin{enumerate}[itemsep=2pt,topsep=2pt]
\item \textbf{Faithful reproduction and audit} (\Cref{sec:repro}). We rebuild the
pipeline to the episode and iteration, reproduce the paper's headline numbers, and
establish that the benchmark's \emph{coded} success criterion is \degc{20} angular
tolerance (its \degc{5} appears in prose only); we report both throughout.
\item \textbf{The subgoal-scale law} (\Cref{sec:scale}). Retraining the drafter at
$S\in\{5,10,15,25\}$ with the control loop matched to the stride reveals an interior
optimum at $S{=}10$ that beats the paper configuration by $+6.1/+10.1$pp and, on the
paper's own protocols, exceeds its published numbers. A cadence-isolation control
identifies \emph{scale matching}, not re-planning frequency, as the mechanism.
\item \textbf{Reality-verified speculative consumption} (\Cref{sec:spec}). Replacing
DSpark's learned verifier with the achieved environment state yields an
${\sim}11\times$ reduction in diffusion compute versus every-step drafting at matched
scale, at neutral success rate.
\item \textbf{Negative results} (\Cref{sec:neg}, \Cref{sec:regress}). Refinement heads,
blind commitment, learned confidence, and a deterministic regressor all hurt in closed
loop despite matched-or-better \emph{offline} fidelity --- causally implicating the
drafter's conditional \emph{spread} as load-bearing.
\end{enumerate}

% ============================================================
\section{Setting and Faithful Reproduction}
\label{sec:repro}

\paragraph{The FF-JEPA planner.} Observations are encoded by a frozen LeWM JEPA encoder,
$z_t=\mathrm{enc}(o_t)$. A diffusion model $G$ (``GDM'': a 3-block DiT with adaLN-Zero
conditioning, $N{=}3$ future subgoals, $\epsilon$-prediction, trained on
success-filtered expert demonstrations) proposes the next subgoal latent
$\hat z_{sg,m+1}$ from the current encoded observation. CEM then optimizes an action
sequence against a terminal-$L_2$ cost to that subgoal, re-planning every $H$
environment steps. The paper fixes the stride to $H{=}25$.

\paragraph{Reproduction (audit closed).} Our faithful retraining reproduces the paper:
short-horizon $95.31\%$ (paper $96.09\%$; same-set oracle $93.36\%$; $n{=}256$),
long-horizon $t{=}75$ multiseed mean $90.0\%$ (paper $91.80\%$), random-init
$96.5\%/78.1\%$ at \degc{20}/\degc{5} (paper $82.42\%$). Pipeline faithfulness is
over-determined: a from-scratch rebuild of the dense training set reproduces the
original training population to the episode ($9{,}094$) and iteration ($171{,}940$),
and the retrained $S{=}25$ drafter reproduces the reference offline fidelity anchor to
four decimals ($0.1593$ vs.\ $0.1593$ at $m{+}1$).

\paragraph{The \degc{20}/\degc{5} criterion.} FF-JEPA's prose states a \degc{5} angular
tolerance, but the environment's shipped evaluation code enforces \degc{20}; the oracle
(true goal, exact pipeline) reaches only ${\sim}84$--$86\%$ at literal \degc{5} even
under matched-population filtering, making a genuine $96\%$ at \degc{5} for a
\emph{sampled}-subgoal method implausible. We therefore treat \degc{20} as the operative
criterion and \degc{5} as a strict \emph{precision} criterion, reporting both throughout.

\paragraph{Protocols.} Two protocols recur below. \emph{Paper protocol}: $n{=}256$,
paper-matched sampling (reproduction and random-init). \emph{Controlled protocol}:
$n{=}62$ fixed held-out episodes, goal offset $150$, budget $300$ steps, block-only
criterion, CEM seeds 42--45 (curve: 42--49); every arm sees byte-identical episodes and
seeds, so within-table deltas are exact treatment effects. Per-arm
$\mathrm{SE}\approx1.7$pp (4 seeds) / $0.9$pp (8 seeds).

% ============================================================
\section{The Subgoal-Scale Law}
\label{sec:scale}

We retrain the drafter at $S\in\{5,10,15,25\}$ with an identical recipe --- only the
target spacing changes, and pair counts are construction-identical --- and scale the
control loop with the stride (CEM horizon $=$ re-plan cadence $=S$). Offline pre-gates
confirmed fine-stride targets are well-posed (displacement $4.7\times$ the encoder noise
floor at $S{=}5$) and that the fine-scale CEM cost landscape \emph{sharpens} (true-plan
ranking percentile $0.93$ at $S{=}5$ vs.\ $0.82$ at $S{=}25$; cost dynamic range
$3\times$ larger).

\begin{table}[h]
\centering
\caption{Success rate vs.\ subgoal scale (controlled protocol). A clear interior optimum
at $S{=}10$: $+6.1/+10.1$pp over the paper configuration ($+8.7$pp at \degc{5} with 8
seeds, $\mathrm{SE}\approx0.9$pp, positive in every seed). $S{=}10$ pays almost no
penalty between the loose and strict criteria; $S{=}25$ drops ${\sim}5$pp.}
\label{tab:scale}
\begin{tabular}{lcccc}
\toprule
& $S{=}5$ & $S{=}10$ & $S{=}15$ & $S{=}25$ (paper config) \\
\midrule
SR \degc{20} (seeds 42--45)  & 92.3 & \ours{97.2} & 96.8 & 91.1 \\
SR \degc{5} (seeds 42--45)   & 91.9 & \ours{96.4} & 94.4 & 86.3 \\
SR \degc{5} (8 seeds, 42--49)& 92.9 & \ours{96.4} & 94.4 & 87.7 \\
Diffusion calls / episode    & 60   & 30          & 20   & 12   \\
\bottomrule
\end{tabular}
\end{table}

\paragraph{Mechanism: scale matching, not cadence.} A cadence-isolation control ---
re-anchor every 5 steps against \emph{unchanged} 25-step targets --- \emph{collapses} to
$78.6\%/77.8\%$ ($-12.5$pp below baseline): each re-draft re-samples a distant subgoal
from a wide conditional posterior, and CEM chases a churning target. Together with the
monotone commitment losses of \Cref{sec:neg}, the correct statement is: \textbf{closed-loop
success is governed by matching the control loop to the subgoal scale; finer matched
scales win, to an interior optimum} --- not ``re-plan more,'' not ``re-plan less.''

\paragraph{Prediction error is horizon-intrinsic.} Across all four drafters, offline
prediction error collapses onto the \emph{absolute} look-ahead horizon, independent of
block position ($+10$ steps: $0.105$ vs.\ $0.112$; $+15$: $0.132$ vs.\ $0.140$), and the
$S{=}5$ drafter's \emph{third} block position ($+15$, $0.140$) outpredicts the $S{=}25$
drafter's \emph{first} ($+25$, $0.159$). ``Suffix decay'' is task-future uncertainty, not
an architectural defect --- retiring intra-block architecture work.

% ============================================================
\section{Negative Results: The DSpark Port}
\label{sec:neg}

We implemented the published DSpark speculative-decoding composition on the GDM and
evaluated each component in the closed loop (controlled protocol, $S{=}25$).

\begin{table}[h]
\centering
\caption{DSpark components ported literally. Lower latent error does not mean better
control: the refiner \emph{halved} offline prediction error while losing 11pp SR ---
offline fidelity anti-correlates with closed-loop SR. The mechanism is conditional-spread
collapse: pulling drafts toward the conditional mean produces less-reachable targets.
This closes refinement, learned-confidence scheduling, and blind commitment as directions.}
\label{tab:dspark-neg}
\begin{tabular}{lccl}
\toprule
Arm & SR \degc{20} & SR \degc{5} & Verdict \\
\midrule
GDM, re-draft every step (reference)          & 91.9 & 83.1 & --- \\
+ refinement head (fixed-1, same cadence)     & 81.1 & 72.2 & refinement $-11$pp \\
+ blind commit-3                              & 73.8 & 50.0 & commitment $-18$pp \\
+ learned-confidence adaptive commit          & 64.5 & 52.0 & worse than blind \\
raw open-loop chain (commit-6)                & 72.6 & 63.7 & healthy chain still loses \\
\bottomrule
\end{tabular}
\end{table}

% ============================================================
\section{Random Initialization: The Paper's Hardest Setting}
\label{sec:randinit}

\begin{table}[h]
\centering
\caption{Random initialization, per-episode random goals (paper III-A protocol). The
protocol bug that once produced a spurious $100\%$ was fixed \emph{before} these runs;
the reproduced baseline scoring $78.1\%$ on the identical code path is the control.
$S{=}10$: \ours{$1024/1024$ episodes over four independent seeds, zero failures at both
tolerances}. The paper's criterion mapping is ambiguous for this cell (its $82.42$ aligns
with our strict-criterion baseline); under either reading, fine scale moves it to $100\%$.}
\label{tab:randinit}
\begin{tabular}{lcc}
\toprule
& SR \degc{20} & SR \degc{5} \\
\midrule
LeWM (flat, goal image; paper)              & \multicolumn{2}{c}{0.00} \\
FF-JEPA Det (paper)                         & \multicolumn{2}{c}{81.25} \\
FF-JEPA GDM (paper)                         & \multicolumn{2}{c}{82.42} \\
FF-JEPA GDM (our faithful repro, $n{=}256$) & 96.5 & 78.1 \\
Ours, $S{=}5$ ($n{=}256\times2$ seeds)      & 100.0 & 94.5 / 99.2 \\
Ours, $S{=}10$ ($n{=}256\times4$ seeds)     & \ours{100.0} & \ours{100.0} \\
\bottomrule
\end{tabular}
\end{table}

% ============================================================
\section{Reality-Verified Speculative Consumption}
\label{sec:spec}

The surviving DSpark idea is its \emph{skeleton}: draft a block cheaply, verify, accept
the surviving prefix. Our port replaces the learned verifier with reality: draft the
$N$-block once; after each hop, compare the \emph{achieved} latent to the subgoal just
pursued; within tolerance $\tau$ $\Rightarrow$ consume the next drafted position (no
diffusion call); otherwise $\Rightarrow$ re-draft from the achieved state. Re-anchoring is
never sacrificed when reality diverges. Offline calibration showed 8-step DDIM drafts are
indistinguishable from 50-step drafts (matched-noise distance $0.03$; spread preserved);
both transferred to the closed loop.

\begin{table}[h]
\centering
\caption{Speculative consumption (controlled protocol). NFE $=$ denoising function
evaluations per episode (calls $\times$ DDIM steps). \textbf{Cost ratios, denominators
pinned:} the headline ${\sim}11\times$ is $1500{\to}135$ against \emph{our own}
every-step $S{=}10$, $k{=}50$ reference (identical scale and loop --- the
apples-to-apples comparison); the $4.4\times$ in this caption is $600{\to}135$ against
FF-JEPA's \emph{native} configuration ($S{=}25$ every-step), which is cheaper per
episode but weaker. The reduction compounds two factors --- fewer calls (ratio
${\sim}0.56$) and cheaper calls ($k{=}8$ vs.\ $50$); the $k$-isolation row holds $k$
fixed to isolate the former; the mirror arms (every-step at cheap $k$, isolating the
consumption policy against equally-cheap-per-call baselines) are reported in the
$k$-curve paragraph below.
At 8-seed power, verified coasting
is \degc{20}-neutral at $S{=}10$ but costs ${\sim}3.5$pp at \degc{5}; at $S{=}5$ the sign
\emph{reverses} (spec-accept beats its own every-step reference on both criteria) --- when
re-drafts churn a close target, skipping verified re-drafts is better control. The full
$S{=}10$ system stays $+6.3/+5.2$pp ahead of the paper configuration (8 seeds) at
$4.4\times$ lower diffusion cost.}
\label{tab:spec}
\begin{tabular}{lcccc}
\toprule
Arm ($S{=}10$ unless noted) & SR \degc{20} & SR \degc{5} & call ratio & NFE/ep \\
\midrule
every-step drafting, $k{=}50$ (reference)          & 97.2 & 96.4 & 1.00 & 1500 \\
spec-accept $\tau{=}0.15$, $k{=}8$                  & 98.0 & 95.2 & 0.67 & 161 \\
spec-accept $\tau{=}0.20$, $k{=}8$                  & \ours{98.0} & \ours{94.8} & \ours{0.56} & \ours{135} \\
spec-accept $\tau{=}0.30$, $k{=}8$                  & 96.8 & 94.8 & 0.45 & 108 \\
spec-accept $\tau{=}0.20$, $k{=}50$ (k-isolation)   & 98.0 & 95.6 & 0.55 & 825 \\
blind commit-2 (no verification)                   & 98.8 & 96.4 & 0.50 & 750 \\
\midrule
$S{=}5$ every-step (reference)                     & 92.3 & 91.9 & 1.00 & 3000 \\
$S{=}5$ spec-accept $\tau{=}0.20$, $k{=}8$          & \ours{98.0} & 94.8 & 0.42 & 202 \\
\midrule
\multicolumn{5}{l}{\emph{8-seed extension (seeds 42--49; $\mathrm{SE}\approx0.9$pp)}} \\
$S{=}10$ every-step                                 & 98.0 & \ours{96.4} & 1.00 & 1500 \\
$S{=}10$ spec-accept $\tau{=}0.20$, $k{=}8$         & 97.4 & 92.9 & 0.56 & \ours{135} \\
$S{=}5$ every-step                                  & 92.9 & 92.9 & 1.00 & 3000 \\
$S{=}5$ spec-accept $\tau{=}0.20$, $k{=}8$          & \ours{98.4} & \ours{96.2} & 0.42 & 202 \\
\midrule
paper config ($S{=}25$ every-step, 8 seeds \degc{5})& 91.1 & 87.7 & --- & 600 \\
\bottomrule
\end{tabular}
\end{table}

\paragraph{Hyperparameter selection.} $\tau$ and $k$ were chosen by offline calibration
and the closed-loop sweep of \Cref{tab:spec}, which runs on the same fixed evaluation
population as the reported cells --- a selection-on-evaluation-set risk we address two
ways. First, $\tau$ was re-selected on an episode-disjoint validation population: the
grid is flat ($63.3$--$64.9$ \degc{20} across $\tau\in[0.15,0.30]$ on the harder
unfiltered population) and $\tau{=}0.20$ remains the selected cell --- the threshold is
not load-bearing. Second, and independently, the same $\tau{=}0.20$ transferred to
Reacher with \emph{zero re-tuning} and remained SR-neutral at call ratios
$0.46$--$0.88$ in every cell of every regime (\Cref{sec:reacher}) --- the behavior of a
robust operating point, not an overfit one.

\paragraph{The $k$-curve, and an honest decomposition of the cost claim.} On the
re-derived pipeline we swept $k\in\{1,2,4,8,50\}$ at matched long-horizon anchors on
both environments. PushT: every-step SR is flat across $k\in\{4,8,50\}$ and cliffs
below ($55$--$56$ at $k\leq2$ vs.\ $63$--$65$ at $k\geq4$); Reacher tolerates even
$k{=}1$ ($95.1$). We therefore \textbf{freeze $k{=}4$ globally} --- the knee of the
cost-quality curve on the stricter environment. This resizes the efficiency claim into
two separable factors: (i) \emph{cheap drafts}, $k{=}50\to4$, a $12.5\times$ per-call
reduction available to every-step drafting too --- FF-JEPA's $k{=}50$ was image-domain
convention, not necessity; and (ii) \emph{verified consumption}, spec-accept's own
marginal factor of ${\sim}1.8\times$ at matched $k$ (call ratio ${\approx}0.55$),
SR-neutral or better. The legacy single-number comparison ($135$ vs.\ $1500$ NFE/ep,
${\sim}11\times$) conflates the two and we retain it only against the paper's native
configuration. The verifier also passes a self-test at the cliff: at $k\leq2$, where
drafts are genuinely poor, the reality check rejects essentially every draft (call
ratio $1.00$) --- it degrades to every-step rather than coasting on bad plans.
Wall-clock, measured with CUDA-synced instrumentation: the drafter is ${\sim}1\%$ of
episode time in this stack (CEM dominates), so the NFE reduction is a drafter-side
compute claim --- material when the drafter is large relative to the planner, immaterial
to end-to-end latency here; fine stride itself costs nothing ($S{=}10$ episodes run
$12\%$ \emph{faster} than $S{=}25$ at matched protocol).

% ============================================================
\section{Is the Diffusion Drafter Necessary?}
\label{sec:regress}

A deterministic MLP regressor ($z_{cond}\to$ block, the offline-strong JointMLP class,
weight decay $10^{-3}$) as the subgoal source, re-anchoring every replan, matches the
diffusion drafter's \emph{offline} prediction error on the same windows ($0.30$ vs.\
$0.32$ at $S{=}10$). Closed loop, it loses everywhere.

\begin{table}[h]
\centering
\caption{Diffusion vs.\ deterministic regression as the subgoal source (seeds 42--45).
\textbf{The drafter's stochasticity is load-bearing}: at matched-or-better offline
fidelity, the zero-spread mean target loses everywhere, and its $-11.7$pp strict-criterion
loss at $S{=}25$ quantitatively mirrors the refinement head's $-11$pp (\Cref{tab:dspark-neg}),
closing the spread thesis causally --- any operator that regresses subgoals toward the
conditional mean damages control.}
\label{tab:regress}
\begin{tabular}{lcccc}
\toprule
& \multicolumn{2}{c}{$S{=}10$} & \multicolumn{2}{c}{$S{=}25$} \\
& \degc{20} & \degc{5} & \degc{20} & \degc{5} \\
\midrule
diffusion drafter (reference) & 97.2 & 96.4 & 91.1 & 86.3 \\
deterministic regressor       & 91.9 & 91.9 & 86.3 & 74.6 \\
\midrule
$\Delta$                      & $-5.2$ & $-4.4$ & $-4.8$ & $-11.7$ \\
\bottomrule
\end{tabular}
\end{table}

\paragraph{Dose-response: spread as a continuous lever.} The regressor, the refinement
head, and blind commitment are all \emph{categorical} spread reductions. To close the
thesis causally we sweep a single continuous lever: the sampler's injected-noise scale
$\gamma$ ($\gamma{=}0$ deterministic, $1$ native, $>1$ inflated), inference-only, on the
fixed long-horizon population ($t{=}100$, seeds 42--45 at the endpoints). SR is
\textbf{non-monotone with an interior optimum at native spread}: for $S{=}25$,
$49.0 \to 57.9 \to 28.7 \to 4.7$ across $\gamma{=}0 \to 1 \to 1.5 \to 2$ (\degc{20});
$S{=}10$ traces the same shape with a gentler low side ($56.6$ at $\gamma{=}0$,
$59.1$--$59.4$ at $\gamma{=}1$--$1.25$, collapse by $\gamma{=}2$). Too little spread
re-creates the regressor's loss; too much destroys the draft distribution entirely.
The three categorical ablations and this curve are one phenomenon: closed-loop control
pays for conditional spread near the level the drafter was trained at, in both
directions.

% ============================================================
\section{Failure Anatomy}
\label{sec:failure}

Instrumented traces (every re-plan's achieved latent and drafted subgoal) over all 8
seeds of $S{=}10$: \textbf{18 failures in 496 episodes, all tolerance-boundary
near-misses} --- final latent distance-to-goal within the success distribution's 95th
percentile (one failure sits \emph{closer} than the median success: a pure angle miss).
Zero unreachable-subgoal failures, zero divergence, zero stuck episodes; successful
episodes' commanded-to-achieved step ratio has median $1.06$. FF-JEPA's own reported
failure mode (``subgoal looks right, agent cannot reach it'') is \emph{eliminated} at fine
scale; the residue is quantization at the success threshold.

% ============================================================
\section{Horizon Scaling (Fixed Population)}
\label{sec:horizon}

Goal offsets $t\in\{25,50,75,100,150\}$ on the \emph{identical} 256 episodes at every
offset (reindexed from the $t{=}150$ set, success-population eligibility), budget $2t$,
seeds 42/43 --- so horizon is the only variable, free of the eligibility confound that
invalidates independent per-offset draws.

\begin{table}[h]
\centering
\small
\setlength{\tabcolsep}{4.5pt}
\caption{Fixed-population horizon sweep ($n{=}512$/cell), re-derived end-to-end from
public assets on independent hardware (anchor cells match the original pipeline within
noise), now \emph{including the flat-planner arm on the identical population} and
extended to $t{=}150$. \textbf{Flat planning collapses ($98.1 \to 2.7$) while $S{=}10$
is flat at ${\sim}98$ across the entire sweep.} Spec-accept tracks every-step within
noise at every offset (call ratio $0.54$ at long horizons, $0.85$ at $t{=}25$: the cost
win is a long-horizon phenomenon). Replan-rate controls (RH$\,{=}1/2$) do not rescue
the flat planner (\S\ref{sec:horizon} text).}
\label{tab:horizon}
\begin{tabular}{lcccccccc}
\toprule
& \multicolumn{2}{c}{LeWM flat} & \multicolumn{2}{c}{$S{=}25$ (FF-JEPA)}
& \multicolumn{2}{c}{$S{=}10$ every-step} & \multicolumn{2}{c}{$S{=}10$ $+$ spec-accept} \\
\cmidrule(lr){2-3}\cmidrule(lr){4-5}\cmidrule(lr){6-7}\cmidrule(lr){8-9}
$t$ & \degc{20} & \degc{5} & \degc{20} & \degc{5} & \degc{20} & \degc{5} & \degc{20} & \degc{5} \\
\midrule
25  & 98.1 & 95.3 & 99.0 & 93.2 & \ours{99.6} & \ours{97.5} & 99.6 & 96.7 \\
50  & 60.4 & 48.8 & 94.1 & 85.0 & \ours{99.0} & 93.8 & 97.1 & 93.0 \\
75  & 42.6 & 28.5 & 95.9 & 82.2 & \ours{98.1} & \ours{95.1} & 97.5 & 93.2 \\
100 & 12.7 & 5.3  & 95.7 & 86.1 & 98.1 & 94.3 & \ours{99.2} & \ours{95.7} \\
150 & 2.7  & 2.0  & 93.9 & 85.6 & 98.4 & 96.5 & \ours{98.8} & \ours{96.9} \\
\midrule
\degc{20} gap ($S{=}10-$flat) & & & & & \multicolumn{2}{c}{$+1.5\to+38.6\to+55.5\to+85.4\to+95.7$} & & \\
\bottomrule
\end{tabular}
\end{table}

\paragraph{Cross-paper regime (VLWM fixed-budget-50).} We additionally evaluate under the
VLWM \citep{vlwm} fixed-budget-50 regime --- a harsher protocol than FF-JEPA's
final-window budget of $2t$, holding the step budget at 50 for \emph{every} offset
(\Cref{tab:vlwm}).

\begin{table}[h]
\centering
\caption{Cross-paper cell: our three arms under VLWM's fixed-budget-50 regime, all
\degc{20}, seed means. Rightmost columns are the published PushT numbers of VLWM and of
LeWM (flat) under that regime. Budget starvation reproduces the published collapse, yet
within the identical budget our arms sit ${\sim}3\times$ above VLWM's bars. Curve
\emph{shapes} are comparable across papers; absolute numbers are not (start/goal sampling
differs), so we cite this as ``same budget regime,'' not a like-for-like win.
$^{\dagger}$At $t{=}25$ the fixed-50 budget equals FF-JEPA's native $2t$ budget, so the two
protocols coincide (the row is the $t{=}25$ cell of \Cref{tab:horizon}); the budget binds
only for $t\geq50$, which is where the published VLWM/LeWM bars begin.}
\label{tab:vlwm}
\begin{tabular}{lccc cc}
\toprule
& \multicolumn{3}{c}{Ours (fixed-budget-50)} & \multicolumn{2}{c}{Published} \\
\cmidrule(lr){2-4}\cmidrule(lr){5-6}
$t$ & $S{=}25$ & $S{=}10$ & $S{=}10$ $+$ spec & VLWM & LeWM (flat) \\
\midrule
25$^{\dagger}$ & 98.1 & \ours{99.8} & 99.6 & 94 & 90 \\
50  & 83.6 & 88.7 & \ours{90.0} & 60 & 46 \\
75  & 55.7 & 59.8 & \ours{60.9} & 20 & 12 \\
100 & 23.4 & 24.4 & \ours{27.3} & 12 & 8  \\
\bottomrule
\end{tabular}
\end{table}

Four read-offs, all consistent with the main story: (a) \textbf{budget starvation
reproduces the collapse} --- confirming VLWM's PushT cliff is largely a budget artifact,
which validates our protocol note; (b) \textbf{within their regime we sit ${\sim}3\times$
above their bars} ($60.9$ vs.\ $20$ at $t{=}75$), citable as ``same budget regime'' with
the start/goal-sampling caveat; (c) \textbf{spec-accept beats every-step at every offset}
under time pressure --- the anti-churn effect (\Cref{sec:spec}) generalizes to $S{=}10$
once budget binds, a new data point for the $\tau$-trade; (d) \textbf{the $S{=}10{-}S{=}25$
gap compresses under starvation} --- scale matching buys precision, not raw speed.

\paragraph{Strongest-baseline control (replan rate).} The flat baseline's replan rate is
a hidden hyperparameter (\Cref{sec:reacher} develops this fully). On PushT it does not
change any comparison above: at the short protocol, faster replanning ties the native
configuration (RH${=}2$: $93.4$ vs.\ RH${=}5$: $93.4$; RH${=}1$: $81.8$; $n{=}1024$
over 4 seeds), so the tables already anchor to the strongest flat configuration. On the
fixed-population sweep, \emph{no} replan rate rescues flat planning from the horizon
collapse --- RH${=}1$ falls $88.3\to18.0$ and RH${=}2$ $96.3\to47.1$ over
$t{=}25\to75/100$, alongside the native configuration's published-anchored cliff ---
while the drafted-subgoal arms of \Cref{tab:horizon} hold ${\geq}97$ throughout.
The in-population flat arm of \Cref{tab:horizon} (RH${=}5$, the strongest
configuration at distance) completes the picture: $98.1\to2.7$ across the sweep.
Replanning faster is not a substitute for planning \emph{toward something nearer}.

% ============================================================
\section{Headline Comparison}
\label{sec:headline}

\begin{table}[h]
\centering
\small
\setlength{\tabcolsep}{5pt}
\caption{Where we stand, both criteria per cell (\degc{20} $=$ the criterion enforced by
the benchmark's shipped code; \degc{5} $=$ the strict precision criterion its prose
states). ``Ours'' modifies FF-JEPA's GDM system only in subgoal scale, loop matching, and
consumption policy --- encoder, drafter architecture, training recipe, and CEM are
unchanged. $^{p}$The paper publishes one number per cell without stating its criterion;
each is placed on the side it empirically aligns with, with $\cdot$ marking the untested
side. Comparisons are valid only within a side.}
\label{tab:headline}
\begin{tabular}{lccccc}
\toprule
& & \multicolumn{2}{c}{FF-JEPA GDM} & \multicolumn{2}{c}{\ours{Ours}} \\
\cmidrule(lr){3-4}\cmidrule(lr){5-6}
Setting (SR \degc{20} / \degc{5}) & LeWM (flat) & published$^{p}$ & our repro & $S{=}10$ & $+$ spec-accept \\
\midrule
Short horizon, $t{=}25$          & 94.5 & 96.1 / $\cdot$ & 95.3 / $\cdot$ & \ours{98.4} / 94.1 & 98.4 / \ours{95.3} \\
Long horizon, $t{=}75$           & 3.5  & 91.8 / $\cdot$ & 90.0 / 78.8 & 97.1 / \ours{93.8} & \ours{97.9} / 91.6 \\
Long horizon, $t{=}150$$^{a}$    & ---  & --- & 91.1 / 87.7 & \ours{98.0 / 96.4} & 97.4 / 92.9 \\
Horizon sweep, $t{=}100$$^{b}$   & ---  & --- & 93.8 / 81.8 & 97.9 / \ours{93.8} & \ours{98.6} / 92.6 \\
Random initialization$^{c}$      & 0.0  & $\cdot$ / 82.4 & 96.5 / 78.1 & \ours{100.0 / 100.0} & --- \\
\midrule
Diffusion NFE / episode$^{a}$    & 0    & --- & 600 & 1500 & \ours{135} \\
\bottomrule
\end{tabular}

\vspace{1mm}
{\footnotesize $^{a}$Controlled protocol, 8 seeds. $^{b}$Fixed-population sweep,
$n{=}512$/cell. $^{c}$$n{=}1024$ over 4 seeds, zero failures at both tolerances;
spec-accept not run in this mode (the drafter is the tested component).}
\end{table}

% ============================================================
\section{Generalization to Further Environments}
\label{sec:generalize}
% All results above are on PushT. The following environments extend the
% subgoal-scale + spec-accept claims beyond a single task. Tables are
% placeholders until the runs land.

Every result above is on PushT. To establish which findings are properties of
hierarchical latent planning rather than of one task, we extend to environments of
increasing structural distance from PushT. Reacher is complete; it cleanly separates
the claims: \textbf{the consumption mechanism transfers unconditionally, while the value
of drafting itself is regime-dependent} --- and the regime boundary is itself the
transferable finding.

\subsection{Reacher}
\label{sec:reacher}

\paragraph{Setup.} DeepMind Control Reacher under the LeWM protocol: frozen LeWM
encoder, success $=$ all joints within $0.05$\,rad of the goal configuration, goal
sampled $t$ steps ahead in a held-out episode (episodes ${\geq}8000$ of $10{,}000$),
budget $2t$ unless stated, $n{=}128$ episodes per seed, seeds 42--45. Unlike PushT's
expert demonstrations, the Reacher dataset is collected by an exploratory SAC policy
\citep{lewm}: trajectories wander rather than beeline. Goal-free drafting is therefore
ill-posed --- the data's continuation from $z_t$ carries no information about an
arbitrary hindsight goal --- so all drafters here are goal-conditioned (hindsight goals,
gap ${\sim}U[1,100]$), trained with the identical recipe of \Cref{sec:scale}. A
goal-free control drafter quantifies this: $8.6$--$11.3\%$ SR at every horizon ---
statistically the random-policy floor ($\approx10\%$) --- versus $60.6$--$95.1\%$
goal-conditioned. The converse cell closes the $2{\times}2$: on PushT's expert data,
goal-conditioning the drafter changes \emph{nothing} ($98.9$ vs.\ $99.4$ \degc{20}
short; parity at $t{=}150$) --- conditioning is necessary exactly when the data does
not already flow toward goals, and redundant when it does.

\paragraph{Native task: drafting has no headroom, and pays a detour tax.}
At the native protocol ($t{=}25$, budget 50) the scale law's \emph{shape} reproduces
--- the interior optimum at $S{=}10$ beats the native $S{=}25$ --- but the entire
drafter curve sits far below flat planning (\Cref{tab:reacher-native}). Three controls
attribute this. First, an \emph{oracle} arm (ground-truth dataset latents as subgoals)
lands at exact parity with matched-loop flat planning ($84.2\pm1.0$ vs.\ $84.0$,
$n{=}512$): a single short reach inside the planner's lookahead gains nothing from
decomposition, bounding all drafters at parity. Second, a \emph{straight-line} oracle
(waypoints interpolated in latent space toward the goal, re-anchored per replan)
isolates decomposition from data quality: even geometrically perfect waypoints collapse
short-range SR monotonically in hop count ($95.5 \to 70.7 \to 21.9$ for interpolation
fractions $1.0/0.5/0.25$; the $1.0$ arm reproduces flat planning within $0.2$pp, a
harness self-check) --- the tax is intrinsic to making the planner land on intermediate
targets, not an artifact of meandering data. Third, the deficit below the bound is
worsened by the \emph{detour tax}: waypoints
imitate the exploratory data's meanders, and with budget $2t$ the wasted steps are
unaffordable (the tax is largest under the harshest budgets; see below). Spec-accept
remains SR-neutral relative to its every-step reference throughout --- the mechanism
does not create the tax, it faithfully delivers the drafter's quality at $0.59$--$0.88$
of the calls.

\begin{table}[h]
\centering
\small
\setlength{\tabcolsep}{5pt}
\caption{Reacher native protocol ($t{=}25$, budget 50, $n{=}512$/cell over 4 seeds).
The scale law's interior optimum reproduces ($S{=}10$ tops the drafter curve) but flat
planning dominates every drafter arm; even oracle subgoals do not beat the goal image.
Spec-accept tracks its every-step reference at each scale (call ratios $0.59$--$0.88$).}
\label{tab:reacher-native}
\begin{tabular}{lccccc}
\toprule
& \multicolumn{4}{c}{drafted subgoals, by scale} & flat (goal image) \\
\cmidrule(lr){2-5}
SR (\%) & $S{=}5$ & $S{=}10$ & $S{=}15$ & $S{=}25$ & best RH (see text) \\
\midrule
every-step $\gdm$ & 52.5 & \ours{60.6} & 54.5 & 53.9 & \multirow{2}{*}{$95.7\pm1.9$} \\
$+$ spec-accept   & 46.1 & 58.6 & 60.9 & 59.8 & \\
\bottomrule
\end{tabular}
\end{table}

\paragraph{Strongest-baseline control: replan rate is a brittle trade-off.}
Flat planning's replan rate (receding horizon RH, replanning every $5\cdot$RH steps)
turns out to be a hidden hyperparameter with opposite optima by regime: fast replanning
wins the short task and collapses at distance, slow replanning the reverse
(RH${=}1/2/5$ at $t{=}25$: $91.6/95.7/83.2$; at $t{=}100$: $51.2/66.6/83.2$; at
$t{=}150$: $46.1/62.3/76.0$). \textbf{Every comparison below is anchored to the
strongest flat configuration per cell}, and no single flat configuration is strong
everywhere --- which is precisely the trade-off drafted subgoals resolve.

\paragraph{Long horizon: the crossover, against the strongest baseline.}
On a fixed 128-pair population eligible to $t{=}150$ (\Cref{tab:reacher-horizon}), the
best flat configuration degrades monotonically with goal distance while the drafter
arms are flat --- the gap widens from $+11$ to $+19$pp, with spec-accept delivering it
at a $0.62$ call ratio. The environment-level picture mirrors PushT's horizon sweep
(\Cref{sec:horizon}) in attenuated form: a free-space reach never becomes \emph{hard},
merely under-specified by a remote goal image. Two follow-up controls sharpen the
attribution. The straight-line oracle \emph{fails} here ($33/44/55\%$ at interpolation
fractions $0.25/0.5/0.75$, far below the drafter's $94.5$ and below even flat
planning's $83.2$): latent-space geometry is not linear over long ranges, so the
learned draft distribution is \emph{necessary}, not replaceable by interpolation ---
the same object that is pure overhead at short range carries irreplaceable information
at long range. And the scale law \emph{flattens} where drafting wins broadly: at
$t\in\{100,150\}$ all of $S\in\{5,15,25\}$ land within $91$--$96\%$ of $S{=}10$'s
$94.5/95.1$ --- the interior optimum is a short-range phenomenon on Reacher, whereas
PushT retains it at $t{=}150$ ($S{=}10/15$ over $S{=}5/25$ by ${\sim}3$--$4$pp).

\begin{table}[h]
\centering
\small
\setlength{\tabcolsep}{5pt}
\caption{Reacher fixed-population horizon sweep ($n{=}512$/cell over 4 seeds, budget
$2t$; identical 128 episode--start pairs at every offset). Flat planning shown at its
strongest configuration per column (RH${=}5$; faster replan rates are lower at every
$t\geq100$). Under budget starvation (fixed 50 for all $t$) the ordering \emph{inverts}
--- flat $64.8$--$76.6$ vs.\ drafters $30.9$--$51.2$ --- the detour tax under time
pressure; drafting requires budget headroom to pay off.}
\label{tab:reacher-horizon}
\begin{tabular}{lccc}
\toprule
SR (\%) & $t{=}100$ & $t{=}125$ & $t{=}150$ \\
\midrule
flat, strongest per cell            & 83.2 & 81.4 & 76.0 \\
$S{=}10$ every-step $\gdm$          & 94.5 & 93.2 & \ours{95.1} \\
$S{=}10$ $+$ spec-accept (cr $\approx 0.62$) & 94.5 & 93.2 & 93.9 \\
\midrule
gap (best drafter $-$ best flat)    & $+11.3$ & $+11.8$ & $+19.2$ \\
\bottomrule
\end{tabular}
\end{table}

\paragraph{Efficiency frontier at the crossover cell.}
Sweeping the acceptance threshold and draft steps at $t{=}100$
($\tau\in\{0.1,\ldots,0.4\}\times k\in\{4,8,16\}$) yields a clean SR-vs-compute
frontier: $k{=}4$ dominates uniformly across $\tau$, and at the frozen
$\tau{=}0.2,k{=}4$ spec-accept reaches $95.5\%$ ($n{=}4$ seeds) at $2.6$ drafter NFE
per replan. Decomposed per \Cref{sec:spec}: the cheap-draft factor alone lifts
every-step drafting to $96.1$ at $4$ NFE (and Reacher tolerates even $k{=}1$ at
$95.1$), so spec-accept's marginal contribution here is ${\sim}1.5\times$ at SR-parity;
the legacy comparison against the $k{=}50$ reference ($90.8$ at $50$ NFE) compounds
both factors to ${\sim}19\times$ \emph{with} an SR gain, and is quoted only with that
decomposition. SR degrades gracefully with $\tau$, as an acceptance threshold should.

\paragraph{What transfers, what does not.} Transfers: the scale law's interior optimum;
spec-accept's SR-neutrality at large compute reduction, with zero re-tuning of
$\tau,k$; the long-horizon value of drafting against the strongest flat baseline.
Does not transfer: any native-task benefit --- on a task with no decomposition headroom
(oracle-bounded) and non-expert data, drafted waypoints are strictly a cost at short
range. We regard the resulting usage rule --- \emph{draft when the goal outruns the
planning window and the budget affords it; plan flat otherwise} --- as a finding, not a
caveat: one greedy attempt to automate the rule inside the verifier (serving the goal
whenever a waypoint fails a latent goal-progress test) collapsed the long-horizon
advantage to the flat baseline and is reported as a negative result.

\subsection{TwoRoom}
\label{sec:tworoom}
A navigation task with a bottleneck between rooms. We report it as a boundary case
rather than a full battery: its per-episode target is randomized and \emph{not rendered
in the observation}, so goal-free drafting is structurally ill-posed there (the
current state contains no information about the goal), and while goal-conditioning
restored planning it remained well below PushT-level fidelity in preliminary runs. This
is the sharpest instance of the data/observability condition that Reacher exhibits in
milder form (\Cref{sec:reacher}), and we scope the drafting claims accordingly rather
than pursue the environment further.

\subsection{OGBench-Cube}
\label{sec:ogbench}
A manipulation benchmark \citep{ogbench} for the strongest generalization claim: multi-object
state and longer effective horizons, where the ${\sim}11$--$20\times$ diffusion-cost
reduction matters most. The public dataset and pipeline requirements are staged; the
port follows the Reacher template.
\tbdblock{OGBench-Cube: headline SR / NFE vs.\ FF-JEPA baseline. Port in progress.}

% ============================================================
\section{Discussion and Limitations}
\label{sec:discussion}

\paragraph{What the story is.} Closed-loop success in this planner is governed by matching
the control loop to the subgoal scale, and the drafter's conditional \emph{spread} is
load-bearing: every operator we tried that collapses drafts toward the conditional mean
(refinement, blind commitment, deterministic regression) damaged control despite equal or
better offline fidelity. The efficiency win is orthogonal --- reality is a free verifier.

\paragraph{Limitations.} (i) FF-JEPA is a preliminary paper with no released code; our
``faithful'' configuration is a principled reconstruction where the paper is silent, not a
verified match on every axis. (ii) Cross-environment generality is now measured rather
than assumed, and it is \emph{conditional}: the consumption mechanism (spec-accept)
transferred to Reacher unchanged and un-retuned, but the value of drafting itself
requires (a) goals beyond the planning window, (b) budget headroom for imperfect
waypoints, and (c) offline data whose trajectories are goal-directed enough to imitate
--- exploratory data forces goal-conditioned drafting (Reacher) and unobservable goals
defeat it (TwoRoom). OGBench-Cube remains \TBD. (iii) The \degc{20}/\degc{5} criterion
mapping for the paper's own cells is inferred from empirical alignment, not stated by
the paper; we mark every cross-paper comparison with the side it is valid on. (iv) The
$t{=}150$ cell is a beyond-paper probe, not a reproduction claim. (v) All PushT results
were re-derived end-to-end from public assets on independent hardware: the closed-loop
anchor cells reproduce (baseline $94.6$ / oracle $92.6$ / $S{=}25$ $94.8$ / $S{=}10$
$99.4$ \degc{20} against $93.4/93.4/95.3/98.4$), certifying the pipeline; offline
drafter diagnostics are \emph{not} comparable across dense-file rebuilds and closed-loop
anchors are the arbiter. (vi) Horizon-sweep populations are eligibility-defined;
success-filtered and unfiltered populations give the same qualitative story (collapse
vs.\ flat) at different absolute levels, are reported separately, and are never
spliced.

% ============================================================
\section{Conclusion}
\label{sec:conclusion}
Subgoal scale, matched to the control loop, is the dominant unexplored lever in FF-JEPA-style
hierarchical latent planning: a fine matched stride ($S{=}10$) beats the paper configuration
across every setting we measured and, on the paper's own protocols, exceeds its published
numbers. Reality-verified speculative consumption recovers ${\sim}11\times$
(PushT) to ${\sim}20\times$ (Reacher, at the frontier's best cell) of the added
diffusion cost at neutral-or-better success rate, transferring across environments and
across goal-free/goal-conditioned drafters with zero re-tuning. On Reacher the
generalization is sharpened into an operating regime: drafted subgoals dominate the
strongest flat baseline exactly where the goal outruns the planning window and the
budget affords imperfect waypoints, and are a measurable cost elsewhere --- while flat
planning's replan rate proves a brittle short-vs-long trade-off on both environments
that subgoals, not replanning frequency, resolve. OGBench-Cube and the wall-clock
efficiency figure are the remaining work.

% ============================================================
\bibliography{references}
\bibliographystyle{iclr2026_conference}

% ---- placeholder bib keys used above (fill references.bib) ----
% \citep{ffjepa} \citep{lewm} \citep{cem} \citep{vlwm} \citep{ogbench}

\end{document}
